\chapter{Enumeration}
\label{chap:enumeration}

% Closed question: What is the smallest admissible mark a ledger can carry?

\begin{chapteressay}

When a Lean program is compiled, the first thing the elaborator does is
build a table of names. Every definition has a name. Every theorem has a
name. Every typeclass instance, every constructor, every implicit argument
carries a name through the elaboration pipeline. The names are the
compiler's primary handle on the program. Without names, the compiler
cannot refer to anything; with names, the compiler can refer to anything
once. The act of naming is the act of admitting something into the
program.

This is the first observation of the book. Before a compiler can compare
two terms, it must be able to refer to them. Before it can refer to them,
it must have named them. Before it can name them, it must have produced
them in some order that lets each one receive a unique designation. The
production of an ordered list of names is enumeration.

Enumeration is older than the compiler. It is older than the digital
computer. Bookkeepers enumerate transactions. Librarians enumerate books.
Astronomers enumerate stars. The act is the same in every case: take an
unordered collection of items, walk through them in some chosen order, and
assign each one a position. The position is the index. The item at
position $i$ is the value. The pairing of $(i, \text{value})$ is the
record entry.

A compiler does the same thing with definitions. Open a fresh Lean file
and write \texttt{def x := 3}. The compiler creates an entry in the
namespace: the symbol \texttt{x} now refers to the value \texttt{3} of
type \texttt{Nat}. Write \texttt{def y := 4}. A second entry is added.
Write \texttt{\#check x + y}. The compiler looks up both symbols in the
namespace, retrieves their values, performs the addition, and reports the
type of the result. Every step is an enumeration: the namespace is the
ledger, the symbols are the indices, and the values are the data the
indices carry.

The compiler's enumeration is not abstract. It is implemented by a
specific data structure --- a hash table, or a B-tree, or a trie --- that
supports two operations: insertion (add a new (symbol, value) pair) and
lookup (given a symbol, return its value). The insertion operation must be
deterministic: the same symbol inserted twice must result in either a
conflict (the symbol is already taken) or an overwrite (the new value
replaces the old, with a record of the change). The lookup operation must
be repeatable: looking up the same symbol twice in succession must return
the same value, unless the namespace has been modified in between.

These requirements are not optional. A compiler that occasionally returned
different values for the same symbol would be useless. A compiler that
silently overwrote symbols would produce incomprehensible programs. The
discipline of enumeration is the discipline of admissibility: the compiler
admits a new symbol only when it can be guaranteed to behave consistently
with all previously admitted symbols. The check is at the point of
insertion, not at the point of use.

The same discipline applies to records of physical measurement. A
thermometer that produces the same reading for different temperatures
would be useless. A scale that produced different weights for the same
object would be untrustworthy. The discipline of an instrument is the
discipline of the compiler: each new mark is admitted only after it is
certified to be consistent with the prior record. The certification is the
instrument's commitment to admissibility.

This is the chapter's first claim. Enumeration is not merely a convenient
way to organize records. It is the algorithmic structure forced by the
requirement that records be admissible. A collection of unordered marks
is not yet a record; it is a heap. To become a record, the marks must be
indexed, the indices must be unique, and the lookup must be repeatable.
The structure imposed by these requirements is exactly an enumeration in
the algorithmic sense.

The second example to make this concrete is the warehouse. A receiving
dock at a distribution warehouse processes pallets of goods. Each pallet
arrives with a manifest, is scanned by a barcode reader, and is
assigned a shelf location in the warehouse. The scanning operation is
the enumeration step: it turns a physical pallet (an unindexed object)
into an indexed record (the barcode is the index, the manifest contents
are the value). After scanning, the pallet exists twice: physically in
the warehouse and informationally in the inventory database. Both
records are necessary. The physical pallet is what the customer will
eventually receive. The informational record is what the warehouse uses
to find it.

If the scanner fails or is bypassed, the pallet enters the warehouse
unindexed. The system does not know it is there. The pallet occupies
shelf space, blocks aisles, and consumes inventory that the system
cannot account for. The disconnect between physical and informational
reality grows over time, and at some point an audit must reconcile the
two by physically walking the warehouse and re-scanning every pallet.
The cost of the audit is approximately the cost of building the index
from scratch.

The warehouse and the compiler are the same machine viewed through
different metaphors. Both have a primary record (the inventory database,
the namespace). Both admit new entries through an indexed scan
(barcode reader, name resolver). Both rely on \texttt{DISTINGUISHABLE}
to detect duplicates. Both rely on $O(\log n)$ indexing to make lookup
fast. Both have a halting boundary: pallets whose barcodes cannot be
read, terms whose elaboration does not terminate.

The closed question is:

\begin{center}
\emph{What is the smallest admissible mark a ledger can carry?}
\end{center}

The answer the chapter develops is the \texttt{CarrierProcess}: a
record consisting of a symbol (the name) and a value (the fact), where
the value's type admits a decision procedure for equality
(\texttt{DISTINGUISHABLE}). Below this minimum, the record is not yet
admissible. Above it, more structure can be layered: ordering (Section
2, B-trees), self-reference (Section 3, Godel numbering), and the
boundary of what can be enumerated at all (Section 4, the Halting
Problem).

The chapter walks through five algorithmic forms that each instantiate
this structure: sequential scan (the slowest enumeration, no
preprocessing), index construction (the B-tree, $O(\log n)$ lookup),
Godel numbering (enumeration of syntactic structures, the first
exposure to incompleteness), the halting boundary (the limit of
enumeration, where the algorithm runs out), and \texttt{CarrierProcess}
in Lean (the typeclass that captures the whole pattern). Each section
is a different view of the same minimum admissible mark.

The book's full argument depends on this minimum being correctly
identified. Subsequent chapters will assume that every operation
operates on marks of this form. If the minimum is wrong --- if the
smallest admissible record is actually larger or smaller than
\texttt{CarrierProcess} + \texttt{DISTINGUISHABLE} --- the rest of the
machinery will not compose correctly. Chapter 1 must therefore be read
carefully. The structure introduced here is the only structure the rest
of the book is permitted to assume without further justification.

\end{chapteressay}

\section{Sequential Scan}
\label{se:sequential-scan}

Consider a list of one million integers stored in a Python list and the task of
finding the value 42. The simplest algorithm is sequential scan: start at index
0, read the value, compare it to 42, advance one index, repeat. If the value
appears, return its index. If the scan reaches the end without finding 42,
return a not-found marker.

The algorithm makes at most one million comparisons. Each comparison is a
single admissible event: the program reads one integer, compares it to 42, and
records a Boolean result. The comparison is the smallest unit of work the
program performs. The integer at position $i$ is not modified by the
comparison; it is only inspected. The scan leaves the list exactly as it was.

This is the cleanest possible expression of what a ledger does. The list is the
prior record. The scan is a sequence of admissible events, each one indexed by
$i$, each one producing one Boolean output. The accumulated record of the scan
is the index of the first matching element, or the certificate that no element
matched. Nothing more is carried; nothing less is sufficient.

The cost of sequential scan is $O(n)$ in the size of the list. For one million
integers and an integer comparison costing a few nanoseconds, the scan completes
in milliseconds. This is fast in human terms and slow in algorithmic terms. The
algorithm makes no use of any structure in the list. If the list happens to be
sorted, the scan does not exploit that fact. If the list happens to have an
index, the scan does not consult it. Sequential scan is the algorithm of pure
admissibility: it works for any list, in any order, with no preprocessing.

The Python idiom \texttt{for i, v in enumerate(lst): if v == 42: return i}
captures this exactly. Every value is a mark. Every comparison is a trial.
Every index $i$ is an enumeration of the trial's position. The function
\texttt{enumerate} is named after the mathematical concept of enumeration: it
turns a list (an unordered collection of values) into a sequence of (index,
value) pairs, restoring the ordering that the iterator implicitly carries.

Sequential scan illustrates the smallest admissible mark in two senses. First,
the comparison \texttt{v == 42} produces a single bit of information: either
this element matches or it does not. One bit is the smallest classical unit of
record. Second, the index $i$ is a natural number, which is the smallest
mathematical structure capable of carrying ordered enumeration. The mark
$(i, v)$ is the pair (index, value) --- the irreducible carrier of the
sequential scan's record.

Without the index, the scan could not know where it is. Without the value, the
scan could not know what it is reading. Without the comparison, the scan could
not say whether what it is reading is what it is looking for. The three
components --- index, value, comparison --- are the minimum tuple a sequential
scan requires.

In Lean, this triple appears as the \texttt{CarrierProcess} structure from
\texttt{Episode1.lean}. The structure has two fields: a symbol (the index name)
and a value (the fact). The carrier carries the value with its name. The name
is not the value, and the value is not the name. The structure preserves both,
distinct, in a single record. The decision procedure \texttt{DISTINGUISHABLE}
is the comparison: given two carriers, the typeclass instance answers whether
they are the same.

A sequential scan is the iterated application of \texttt{DISTINGUISHABLE} along
an enumeration of carriers. The Python loop \texttt{for i, v in enumerate(lst)}
is, in Lean terms, the enumeration. The body \texttt{if v == 42} is the
\texttt{DISTINGUISHABLE} application. The return statement is the closure of
the search when a match is found.

The Lean compiler, when it elaborates a sequential scan, performs exactly the
same pattern. Every name in the namespace is enumerated. Every name's type is
compared to the target. The compiler's name resolution is sequential scan over
the namespace environment. The compiler is, at this level, doing what the
program describes. The instrument and the artifact have the same structure.

The lesson of this section is the following: the smallest admissible mark a
ledger can carry is a pair of an index and a value, together with a comparison
that can decide whether two such pairs are the same. Anything less is not yet
a ledger entry. Anything more is the chapter's later sections: indices that
preserve order (next section), encodings that preserve self-reference (Godel),
limits that name the unreachable (halting boundary), and the typeclass that
captures the whole pattern (CarrierProcess in Lean).

The sequential scan is the slowest possible enumeration --- but it is also the
algorithm that requires the least structure. Every faster algorithm in this
chapter (B-tree indices, hash tables, binary search) demands prior structure on
the list. Sequential scan demands nothing. It is the baseline against which
every other algorithm is measured. It is also the algorithm the universe must
support before any other algorithm can be defined.

\section{Index Construction}
\label{se:index-construction}

A database table with ten million rows is not searched by sequential scan. The
scan would visit each row once: at one microsecond per row, the search would
take ten seconds. For an interactive query, ten seconds is too slow. The
solution is the B-tree index.

A B-tree is a balanced search tree whose nodes contain sorted keys and pointers
to child nodes. A typical B-tree has a branching factor of around 100: each
internal node points to roughly 100 children, and the height of the tree grows
as $\log_{100}$ of the number of records. For ten million rows, the tree has
about four levels. A range query that would have required ten million
comparisons under sequential scan requires four pointer dereferences and one
sequential walk through the matching range. The query completes in
microseconds.

Index construction is the act of building this tree. Given an unsorted table,
the database constructs the B-tree by sorting the keys and then organizing
them into the balanced structure. The construction cost is $O(n \log n)$. Once
the tree is built, every query is $O(\log n)$. The construction is performed
once; the queries are performed many times. The index pays for itself after
$n \log n / \log n = n$ queries amortized into the cost.

The B-tree is an enumeration of the records in key order. This is a stronger
enumeration than the one provided by sequential scan: the sequential scan
enumerates in insertion order; the B-tree enumerates in key order. Two records
may have arrived at the database in different orders but receive the same
B-tree index. The index imposes a canonical enumeration on the records.

Two properties of the B-tree matter for the chapter's argument. First, the
index preserves the order of the keys. Given any two records with keys $k_1$
and $k_2$, the index can answer in $O(\log n)$ time whether $k_1 < k_2$,
$k_1 = k_2$, or $k_1 > k_2$. The index carries the ordering relation as a
data structure, not as a runtime comparison: the structure of the tree
encodes the order.

Second, the index supports range queries. A query of the form ``give me all
records with key between $a$ and $b$'' is answered by finding $a$ in the tree
(cost $O(\log n)$), then walking forward through successive keys until $b$ is
exceeded. The walk's cost is $O(k)$ where $k$ is the number of records in the
range. The total cost is $O(\log n + k)$. Sequential scan would require
$O(n)$ regardless of $k$. For $k \ll n$, the B-tree is dramatically faster.

The cost of these capabilities is the storage and maintenance of the tree.
Every insertion, deletion, and update must adjust the tree to maintain its
balance. The amortized cost of an update is $O(\log n)$. For workloads with
many more reads than writes, the cost is acceptable; for workloads with very
high write rates, the index may need to be deferred or recomputed in batches.

The Lean analog is the typeclass \texttt{Number.le} in \texttt{Episode1.lean}.
\texttt{Number.le} is the order relation on natural numbers. The Lean kernel
uses it constantly during elaboration: every comparison of universe levels,
every constraint on numeric literals, every well-founded recursion proof
consults \texttt{Number.le}. The kernel maintains an internal structure that
makes these queries fast. This internal structure is, conceptually, a B-tree
over the kernel's type universe.

\texttt{IndexingProcess} in \texttt{Episode1.lean} captures the abstract
structure of indexing: a process that takes an unordered collection and
produces an enumeration in some canonical order. The typeclass
\texttt{COUNTABLE} certifies that an enumeration exists. Together, they
formalize what a B-tree provides: a countable enumeration of a collection,
indexed in a canonical order, with $O(\log n)$ random access.

The interesting boundary appears when the collection is infinite. A B-tree on
ten million records is finite. A B-tree on the natural numbers is not. The
\texttt{Rational} type in \texttt{Episode1.lean} is countable but infinite.
The \texttt{IndexingProcess} for rationals is the Cantor diagonal enumeration:
list rationals in order of $|p| + |q|$ where the rational is $p/q$. This
enumerates every rational exactly once, in finite time per element, but no
finite stage of the enumeration covers all rationals. The index is countable
but unbounded.

This boundary --- finite versus infinite enumeration --- is the first
appearance of the chapter's halting boundary. A finite enumeration always
terminates. An infinite enumeration terminates only when the query is found,
or not at all if the query is not in the list. The B-tree on ten million
records always answers in $O(\log n)$. The diagonal enumeration of rationals
answers ``yes, $p/q$ is in this list'' in finite time, but answers ``no, the
target is not in this list'' only by exhausting the enumeration, which is
infinite for irrational targets.

The compiler analog of the B-tree is the namespace tree. Lean organizes all
defined names into a tree: \texttt{Nat.add}, \texttt{Nat.mul},
\texttt{List.length}, \texttt{List.append}, and so on. The tree is balanced
implicitly by the namespace structure. When the elaborator looks up the name
\texttt{Nat.add}, it walks the tree from the root, branching at each dot.
The lookup is $O(\log n)$ in the size of the namespace. This is the
compiler's index, and it is built the same way a database B-tree is built.

\section{Godel Numbering}
\label{se:godel-numbering}

Kurt Godel, in 1931, observed that a sentence in a formal language can be
assigned a natural number. The assignment is an effective coding of
syntax by natural numbers: a computable injection from the
sentences of the language into $\mathbb{N}$, with a computable inverse
on its image. Not every natural number is the code of some sentence
(most naturals fail to be valid encodings under the chosen scheme),
so the map is not a bijection onto $\mathbb{N}$; what matters is that
the coding and its decoder are both effective. Once the assignment is
fixed, statements about the language can be expressed as statements
about numbers, and statements about numbers can be encoded back into
statements in the language.

The construction is mechanical. Assign each symbol of the language a small
integer: $\neg \mapsto 1$, $\vee \mapsto 2$, $\forall \mapsto 3$,
$( \mapsto 4$, $) \mapsto 5$, $0 \mapsto 6$, $\textrm{successor} \mapsto 7$,
and so on for every variable and every logical connective. A sentence is a
sequence of symbols, say $s_1, s_2, \ldots, s_n$ with integer codes $c_1,
c_2, \ldots, c_n$. The Godel number of the sentence is
$2^{c_1} \cdot 3^{c_2} \cdot 5^{c_3} \cdot \ldots \cdot p_n^{c_n}$, where
$p_n$ is the $n$-th prime.

Every sentence receives a unique Godel number. By the fundamental theorem of
arithmetic, the prime factorization of an integer is unique, so the Godel
number can be decoded back into its sequence of symbol codes, and from there
(when the codes correspond to a well-formed sentence) into the original
sentence. The encoding is reversible on its image. The encoding is also
computable in both directions: given a sentence, computing its Godel number is
straightforward; given a Godel number, computing its sequence of primes and
their exponents is straightforward (if expensive for large numbers), and
checking whether the decoded sequence is a valid sentence is decidable.

Why does this matter? Because once sentences are numbers, proofs can be
numbers too. A proof is a sequence of sentences, each one either an axiom or
derived from earlier sentences by a fixed set of inference rules. The
sequence of sentences has a Godel number (the product of primes raised to the
sentences' Godel numbers). Every proof is an integer. The set of valid proofs
is a recursively enumerable set of integers. The property ``$p$ is a valid
proof of $q$'' is a decidable relation between two integers $p$ and $q$.

This is the foundational move of Godel's incompleteness theorem. Once
sentences and proofs are numbers, a sentence in the language can be about its
own Godel number. Self-reference becomes admissible. The sentence ``this
sentence has no proof'' is no longer a paradox stated in informal language;
it is a sentence in formal arithmetic whose Godel number $g$ satisfies a
specific arithmetical property: there is no integer $p$ such that $p$ is a
proof of the sentence with Godel number $g$. If the formal system is
consistent, this sentence is unprovable; if the formal system can prove its
own consistency, the sentence is also provable, contradiction.

The boundary Godel numbering exposes is the boundary of formal completeness.
Every consistent formal system rich enough to express arithmetic contains
true statements that the system cannot prove. The system's enumeration of
provable statements is a strict subset of the system's enumeration of true
statements. The two enumerations differ. Godel numbering is the algorithmic
tool that makes the difference visible.

For the chapter's argument, Godel numbering shows that enumeration is not
neutral. Sequential scan enumerates a finite list. Index construction
enumerates a finite or countably infinite collection. Godel numbering
enumerates the syntactic structures of a formal language. Each enumeration
is a different commitment to what can be carried as a mark.

The Lean analog of Godel numbering is the compiler's term representation.
Every Lean term has an internal representation as a tree of constructors:
\texttt{Expr.app}, \texttt{Expr.lam}, \texttt{Expr.forall}, \texttt{Expr.var},
\texttt{Expr.const}, and so on. The tree can be flattened into a sequence of
integers using a similar prime-factorization scheme, or equivalently, into a
byte sequence using a serialization protocol. The compiler can hash this
sequence to produce a unique identifier for the term. Two terms are
definitionally equal exactly when their canonical hash matches. The hash is
the compiler's Godel number for the term.

This matters when the compiler must answer the question ``have I seen this
term before?'' During elaboration, the compiler builds many intermediate
terms, and many of them are repeated. Without a Godel-like indexing, the
compiler would have to compare each new term against every previously seen
term, in $O(n)$ time per query. With the hash, the compiler can index every
seen term in a hash table and answer the query in $O(1)$ expected time. The
elaborator's speed depends critically on this caching, and the caching
depends on the hash, and the hash is a Godel number.

The boundary Godel numbering exposes within the compiler is the same as the
boundary Godel exposed in arithmetic. The compiler can enumerate all
elaboratable terms. It cannot enumerate all type-correct terms, because some
type-correct terms require infinitely many elaboration steps to construct,
and the elaborator must time out. The boundary between elaboratable and
type-correct is the compiler's analog of the boundary between provable and
true. The Godel sentence's analog is a Lean term that is type-correct in
principle but that no finite elaboration trace can reach.

Self-reference in Lean appears as recursive definitions: \texttt{def f : Nat
to Nat := fun n => if n = 0 then 0 else f (n - 1) + 1}. The definition
mentions \texttt{f} inside its own body. The compiler accepts this only after
verifying that the recursion terminates. The termination proof is itself a
Godel-like construction: the proof shows that some well-founded relation
decreases with each recursive call. Without the termination proof, the
self-reference is inadmissible. With it, the self-reference is admitted, and
the function is added to the namespace. The compiler's handling of recursion
is the operational form of Godel's argument: self-reference is admissible
exactly when the formal system can certify its own termination.

\section{The Halting Boundary as Enumeration Limit}
\label{se:halting-boundary}

Type \texttt{\#eval Nat.decEq 3 3} into a Lean REPL. The compiler responds
immediately with \texttt{isTrue rfl}. The expression decides whether 3 and 3
are equal as natural numbers; they are; the proof of equality is the
reflexivity term \texttt{rfl}. The whole computation, including elaboration
and evaluation, takes a few milliseconds.

Now write the following: a program that searches the natural numbers for a
counterexample to Goldbach's conjecture. Goldbach's conjecture says every
even integer greater than 2 is the sum of two primes. The program iterates
$n = 4, 6, 8, 10, \ldots$, and for each $n$, checks whether $n$ can be
written as $p + q$ for some primes $p$ and $q$. If a counterexample is
found, the program halts and returns it. If no counterexample exists, the
program runs forever.

As of the date of this writing, Goldbach's conjecture is still open. No
counterexample has been found. If you run the program, it has been running
for as long as you have been running it. It may halt tomorrow with a
counterexample. It may run for the heat death of the universe without
halting. The mathematics does not currently know which.

The halting boundary is the boundary between programs of the first kind
(\texttt{Nat.decEq 3 3}) and programs of the second kind (the Goldbach
search). Programs of the first kind halt and produce an output. Programs of
the second kind may halt or may not. The boundary is the question: given a
program, does it halt?

Turing proved in 1936 that this question is undecidable. There is no
program $H$ that takes another program $p$ as input and decides whether $p$
halts. The proof is by contradiction: suppose $H$ exists. Construct $D$ as
follows: $D(p)$ runs $H(p)$; if $H(p) = $ halts, then $D(p)$ loops forever;
if $H(p) = $ does not halt, then $D(p)$ halts. Now ask: does $D(D)$ halt?
If $H(D) = $ halts, then $D(D)$ loops forever, so $H(D)$ should have said
``does not halt.'' If $H(D) = $ does not halt, then $D(D)$ halts, so $H(D)$
should have said ``halts.'' Either way, $H$ contradicts itself. Therefore
$H$ does not exist.

The halting boundary is the limit of enumeration. The enumeration of halting
programs is recursively enumerable: a universal Turing machine can simulate
each program in turn and add to the list those that halt. The enumeration of
non-halting programs is not recursively enumerable: there is no algorithm
that can certify, in finite time, that a particular program does not halt.
The two enumerations are not symmetric. The asymmetry is the halting
boundary.

For the chapter's argument, the halting boundary is the smallest example of
a structure that the ledger can name but cannot enumerate. The ledger can
carry the statement ``this program halts.'' The ledger can carry the
statement ``this program does not halt.'' What the ledger cannot do is
decide, in general, which of these statements is true for an arbitrary
program. The decision procedure is not admissible. The decision is in the
domain of the noncomputable.

In Lean, this boundary appears as the difference between \texttt{Decidable}
and \texttt{Prop}. A decidable proposition has a decision procedure: an
algorithm that, given the data, returns either a proof or a disproof in
finite time. The proposition \texttt{3 = 3} is decidable. The proposition
``program $p$ halts'' is, in general, not decidable. The proposition
``Goldbach's conjecture is true'' is also, as far as we currently know, not
decidable, although it may turn out to be either provable or refutable
through some yet-undiscovered argument.

Lean accepts undecidable propositions as propositions. It does not require
that every proposition be decidable. What Lean requires is that the proofs
the user supplies be checkable in finite time. If you assert that some
program halts and you supply a proof (perhaps a termination proof using a
well-founded relation), Lean checks the proof and admits the claim. If you
assert it without proof, you must use \texttt{sorry} to acknowledge the gap.

The \texttt{ChaitinsNumberSequence} typeclass in \texttt{Episode3.lean}
formalizes the halting boundary. The sequence is the binary digits of
Chaitin's constant $\Omega = \sum_{p \text{ halts}} 2^{-|p|}$, the
probability that a randomly chosen program halts. The first $n$ bits of
$\Omega$ encode the halting status of the first $2^n$ programs. The
sequence is well-defined; the value of any particular bit is provably
unknowable without solving the Halting Problem for the corresponding
programs.

Lean's representation of \texttt{ChaitinsNumberSequence} is a type whose
elements name positions in the sequence, not values. The typeclass admits
the question ``what is the $n$-th bit?'' as a well-formed type-level query.
It does not provide an algorithm to answer the query. The query is a
metavariable that the compiler can carry but cannot resolve. This is the
algorithmic boundary at the type level: the language can express the
question; the compiler cannot answer it.

The halting boundary thus appears at three levels in this chapter. At the
level of programs, it is the undecidability of Turing's argument. At the
level of arithmetic, it is the incompleteness of Godel's argument (every
incompleteness theorem is, in the right perspective, a halting argument).
At the level of types, it is the gap between the decidable and the merely
provable. All three are the same boundary, viewed through three different
instruments. The chapter's claim is that the ledger can carry events on
either side of this boundary --- halting programs as marks, non-halting
programs as named-but-unresolved metavariables --- but no admissible
algorithm can produce a finite enumeration that distinguishes the two
classes for all inputs.

\section{CarrierProcess in Lean}
\label{se:carrierprocess-in-lean}

The Lean file \texttt{device/Measurement/Episode1.lean} defines a structure
named \texttt{CarrierProcess}. The definition is intentionally minimal:

\begin{leanexcerpt}{CarrierProcess}
structure CarrierProcess (Carrier : Type) where
  symbol : Fact
  value  : Number
  event  : Number $\to$ Number := ...
\end{leanexcerpt}

A \texttt{CarrierProcess} is parameterized over a \texttt{Carrier} type
and has three fields: a \texttt{symbol} (a \texttt{Fact}, the bookkeeping
that something has been carried), a \texttt{value} (a \texttt{Number}, the
data being carried), and an \texttt{event} function from \texttt{Number}
to \texttt{Number} (the rule for advancing the value as the carrier
progresses). All three slots must be provided to construct one. The
structure is the algorithmic encoding of the chapter's first claim: a
value is never naked; it arrives with a carrier and a rule for how it
will be carried forward.

A reader of a \texttt{CarrierProcess} can extract the symbol, the value,
or the event rule, but not collapse them into a single quantity. The
three fields are distinct slots in the same record. Construction
proceeds by providing the symbol and value and accepting the default
\texttt{event} (or supplying a specific one); if any field has the wrong
type, the elaborator rejects the construction.

The companion typeclass is \texttt{DISTINGUISHABLE}:

\begin{leanexcerpt}{DISTINGUISHABLE}
class DISTINGUISHABLE
    (Value : Type)
    (Observation : CarrierProcess Value)
  where
  fact         : Fact
  symbol       : Value
  distinct?    : Value $\to$ Bool
  different?   : Value $\to$ Value $\to$ Prop
  dec_distinct : DecidablePred distinct?
\end{leanexcerpt}

A type \texttt{Value} together with an observation
\texttt{CarrierProcess Value} has an instance of
\texttt{DISTINGUISHABLE} when the class provides: a \texttt{Fact} that
the distinction has been recorded, a representative \texttt{symbol}, a
unary predicate \texttt{distinct?} on values, a binary predicate
\texttt{different?} for pairs, and a decidability witness
\texttt{dec\_distinct} for the unary predicate. The decidability witness
is the project's analog of an equality decision procedure;
\texttt{DecidablePred} in Lean is the class of predicates with
constructive truth values,
each carrying the relevant proof.

The typeclass is the formal statement of the comparison requirement. A
sequential scan requires \texttt{DISTINGUISHABLE} on the element type:
without it, the scan cannot decide whether the current element matches the
target. A B-tree index requires \texttt{DISTINGUISHABLE} (and an order
relation): without it, the tree cannot decide where a new key belongs. The
Godel numbering requires \texttt{DISTINGUISHABLE} on the encoding alphabet:
without it, decoding cannot recover the original symbol from its integer
code.

The Lean compiler checks for \texttt{DISTINGUISHABLE} instances at
elaboration time. If a user writes a function that requires
\texttt{DISTINGUISHABLE} on some type, and no instance has been provided,
the elaborator rejects the function. The constraint is enforced at compile
time, not at runtime. There is no possibility of a runtime error from a
missing decision procedure; the compiler refuses to admit code that lacks
the certification.

Running \texttt{\#check @DISTINGUISHABLE} in a Lean session displays the
typeclass's signature. The exact signature shown depends on the pinned
toolchain and the surrounding imports; under this project's class
cascade, the signature is approximately
\texttt{DISTINGUISHABLE : (Value : Type) $\to$ CarrierProcess Value $\to$ Type}.
\textbf{TRACE STATUS: pending verification.} The output above is
reconstructed from the source; it has not been captured from a
run under a pinned toolchain. The decision procedure is provided by
the user (or by the standard library) for each concrete type.
\texttt{Nat} has an instance: equality of natural numbers is
decidable. \texttt{String} has an instance: string equality is decidable.
Function types do not have an instance in general: two functions may agree
on every observable input but differ at internal representations the kernel
cannot inspect. The compiler is honest about this: it refuses to certify
function equality without further information.

The pair \texttt{CarrierProcess} and \texttt{DISTINGUISHABLE} captures the
chapter's central claim in code. A mark has a carrier (a symbol naming
it) and a value (the data). To compare two marks, the carrier types must
admit a decision procedure. The compiler enforces both: the structure
construction requires both fields; the typeclass instance is required for
any operation that compares marks. The enforcement is at the elaboration
boundary --- the compiler will not let you compile a program that violates
these constraints.

This is the answer to the chapter's closed question. The smallest
admissible mark a ledger can carry is a \texttt{CarrierProcess}: an indexed
record consisting of a symbol and a value, with the value's type admitting
\texttt{DISTINGUISHABLE} so that two marks can be compared. The structure
is the simplest record from which the rest of the book's machinery is
built. Every subsequent chapter --- comparison, gap, combination, distance,
motion, transport, reference, curvature, invariance, closure --- assumes
that the marks it operates on are \texttt{CarrierProcess} values whose
underlying types are \texttt{DISTINGUISHABLE}.

The remaining typeclasses in \texttt{Episode1.lean} layer additional
structure on top. \texttt{Natural} and \texttt{CountingProcess} together
form the \texttt{ADMISSIBLE} typeclass: a process that produces natural
numbers as marks. \texttt{Rational} and \texttt{IndexingProcess} together
form \texttt{COUNTABLE}: a process whose marks can be enumerated by natural
numbers. \texttt{Sequence} and \texttt{LimitProcess} together form
\texttt{ENCODED}: a process whose marks form a convergent sequence.
\texttt{Limit}, \texttt{CauchyProcess}, and \texttt{Sample} together form
\texttt{RESIDUE}: a process whose limit point may not be in the carrier's
type, leaving a residue that the next refinement must resolve.

These layered typeclasses are the algorithmic vocabulary the book uses.
Each chapter beyond this one assumes that the marks it operates on belong
to some combination of these typeclasses. The chapter's closing claim is
that any data structure rich enough to carry an admissible measurement
record must instantiate at minimum \texttt{CarrierProcess} and
\texttt{DISTINGUISHABLE}. Lighter structures (a bare value, a number
without a name) are not yet records. Heavier structures (sorted indices,
hash tables, B-trees) are admissible compositions of the same primitives.

\begin{bridge}

The chapter's question was what the smallest admissible mark a ledger can
carry is.

The answer is a \texttt{CarrierProcess}: a value with a symbol that names it,
where the value's type admits \texttt{DISTINGUISHABLE} --- a decision
procedure for equality. Sequential scan applies \texttt{DISTINGUISHABLE}
along an enumeration. B-tree indices precompute the order so that the
scan's cost drops from $O(n)$ to $O(\log n)$. Godel numbering shows that
the enumeration extends to self-referential structures, with the cost of
exposing the halting boundary. The halting boundary is the algorithmic
limit of enumeration: programs of the second kind --- the ones that may
not halt --- form a set that no algorithm can completely enumerate.

What remains is that two carriers carrying different values are not yet a
comparison; they are two records. To compare them, the instrument must do
more than admit \texttt{DISTINGUISHABLE} on their types --- it must
compose two decisions into an output that names which is greater, lesser,
or equal, and it must repeat that composition reliably. Chapter 2 asks
how the instrument carries that comparison without pretending the
comparison is the measurement itself.

\end{bridge}

\begin{coda}{The Warehouse Receiving Dock}

The receiving dock is at the back of the building, opposite the customer
entrance. The dock has three loading bays, a conveyor belt, a scanning
station, and a small office for the shift manager. From six in the morning
until ten at night, trucks back into the bays and unload their pallets
onto the belt. The belt feeds the pallets past the scanning station and
into the warehouse proper, where they are routed to shelves by an
automated system.

Today is a Thursday in late October. The dock manager has been at her job
for eleven years. She knows the rhythm of the shifts, the regular drivers,
the seasonal load patterns, and the small failure modes that occur often
enough to require named procedures. She does not know what is on every
truck before it arrives, but she knows that whatever arrives will pass
through the same sequence of operations: unload, scan, sort, route, shelve.

The fourth operation in the sequence --- routing --- is where the dock's
records become a working inventory. Before routing, a pallet is a physical
object with a paper manifest. After routing, the pallet is an entry in the
warehouse management system: a barcode, a SKU, a quantity, a location, and
a timestamp. The transition between the two states is the scan.

A pallet arrives at the scanning station and is read by an overhead
barcode scanner. The scanner emits a beep, an LED flashes green, and the
pallet's barcode becomes a row in the database. The row contains five
fields: the barcode (a unique identifier), the SKU (the product code), the
quantity (a number of units), the location (the shelf address to which
the pallet will be routed), and the timestamp (the moment of the scan).

The barcode is the carrier symbol. Before the scan, the pallet exists
physically but is invisible to the system. After the scan, the pallet has
a name in the database. The name is unique: no two pallets share a
barcode. The barcode is what the system uses to refer to the pallet for
all subsequent operations.

The SKU is the value. It is the data the carrier carries: which product
this pallet contains. The SKU is not unique --- many pallets contain the
same SKU --- but the pairing of barcode (unique) and SKU (shared)
preserves the source identity of every pallet while permitting aggregate
queries (how many pallets of this SKU do we have?).

The scan is the application of \texttt{DISTINGUISHABLE} to the barcode.
The scanner reads the barcode and compares it to every barcode currently
in the database. If a match is found, the pallet is being rescanned ---
the system flags this as a duplicate and routes the pallet to an
exception bay for the manager to inspect. If no match is found, the
barcode is admitted as a new record. The decision is the comparison; the
output of the decision is the admission or the flag.

The dock processes about two thousand pallets a day. Without an index,
the matching step would take seconds per scan: the database would have
to compare the new barcode to each of the warehouse's hundred thousand
existing barcodes. With a B-tree index on the barcode column, the match
takes microseconds. The index was built once, when the warehouse was
commissioned. Every new pallet's barcode is inserted into the index; the
index grows by one entry per scan. The amortized cost of an insertion is
$O(\log n)$. The match cost is also $O(\log n)$. Both are well below the
human-perception threshold.

The location field is the routing decision. The system computes, based
on the SKU and current inventory, the optimal shelf address for the
incoming pallet. The decision is itself a small search: find a shelf that
has space, is in the correct climate zone, is reachable by the forklift,
and is closest to the consuming workstation. The search is done over the
warehouse's location index, which is another B-tree, this one keyed by
shelf address. The two indices --- barcode and location --- coexist in
the same database, providing fast lookup along two different axes.

Self-reference enters the system at the level of the order book. The
warehouse management system maintains a list of outstanding orders: each
order is a request for $n$ units of a particular SKU, to be picked from
the shelves and shipped to a customer. The list of orders refers to the
list of pallets (which SKUs are in stock), and the list of pallets
refers to the list of orders (which orders are waiting on which SKUs).
The two lists are recursively defined. The system handles this by Godel-
numbering each record: every order has an order ID, every pallet has a
barcode, and the cross-references are stored as integer pointers. The
encoding is reversible: given a pointer, the system can dereference it to
the actual record.

Toward the end of the shift, a problem arises. A truck arrives with a
pallet whose barcode is unreadable --- the label has been damaged in
transit. The scanner reads garbage. The system flags the pallet as
unenumerable: it has no admissible mark. The dock manager calls the
exception procedure: the pallet is moved to a holding area, a new
barcode is generated by the system (with the unique-identifier guarantee
preserved), a paper trail is created linking the new barcode to the
truck's manifest, and the pallet is reintroduced to the scanning station
with its new label. The system now admits it. The original barcode is
recorded in the audit log as ``superseded by'' the new one. The audit
log carries the residue.

This is the halting-boundary analog for the dock. Some pallets cannot be
enumerated by the standard procedure. The procedure does not enumerate
them; it does not pretend they do not exist; it transfers them to the
exception process. The exception process is, in algorithmic terms, a
hand-coded resolution that the elaborator (the dock manager) performs
when the standard mechanism cannot proceed. The system's correctness
depends on this exception process being available and reliable.

The far end of the warehouse, behind a roll-up door, is the section
that has not yet been racked. The construction crew finished pouring
the concrete last month, but the shelving has not yet been installed.
Pallets cannot be stored there because the location index has no
addresses for that section. From the warehouse management system's
perspective, the unracked section does not exist: there is no row in
the location table that points to it. Physically, the section is real;
algorithmically, it is invisible. The unracked section is the dock's
analog of the unenumerable: a region of the building that the system
knows about but cannot route to.

At ten in the evening, the last truck leaves. The shift manager closes
the dock door, sets the security alarm, and locks the office. The day's
scan log contains roughly two thousand new entries, one per pallet. Each
entry has the five fields. Each pallet is now \texttt{COUNTABLE} within
the system: it has a barcode, a SKU, a quantity, a location, and a
timestamp. The shelf occupancy data, aggregated across the day, becomes
input to tomorrow's routing algorithm. The system's record has grown by
two thousand marks.

The dock is the most ordinary instance of what the chapter has named.
The barcode is the symbol. The SKU is the value. The scan is the
\texttt{DISTINGUISHABLE} decision. The B-tree index is the
\texttt{IndexingProcess} that makes the decision fast. The unenumerable
pallet is the exception that proves the halting boundary. The unracked
section is the part of the world the algorithm cannot reach. Every
pallet is a \texttt{CarrierProcess}. Every scan is the admission of one
mark into the ledger.

\end{coda}
